\chapter*{Preface}

\addcontentsline{toc}{chapter}{Preface}

This script is the result of my ceaseless, unwavering interest in plasma physics and was originally composed as a personal draft of introductory notes on the subject. Fortunately, the substantial effort invested in the craft of such a script has further motivated me to reformat it into a comprehensive review of plasma physics, which has, before long, become far more detailed than any introductory script in the relevant literature. Each chapter of the script initially addresses its topic at an introductory level, thereby motivating readers to continue following my theoretical development of the subject at an advanced level. The script thus presents a complete graduate-level review of plasma physics.

\quad The first chapter is dedicated to the theoretical development of the underlying concepts in plasma physics, hence the title. We present the subject's main hypotheses and discuss their relevance in particular regimes. The topics covered include Debye shielding and quasi-neutrality, which together describe the response of plasmas to small-amplitude perturbations in the electrostatic potential, as well as transport phenomena at an introductory level. By creating the link between Debye shielding, quasi-neutrality, and collisional processes, we shall establish the fact that plasmas are weakly collisional.

\quad The second chapter offers a complete review of charged particle motion in various configurations of fields, thus exploring the dynamics of plasma particles. The concept of drift velocity proves fundamental in the analytical approach to the kinematics of charged particles. This chapter also introduces readers to non-conservative phenomena and to the concepts of adiabatic invariants and symmetry breaking.

\quad The third chapter elaborates on plasma dynamics from the statistical mechanical point of view by introducing the phase space and distribution functions within it. These considerations will serve as a foundation for the theoretical development of the three main models of plasma dynamics: the Vlasov model, the two-fluid model, and the magnetohydrodynamic model, each of which views plasma behaviour from a perspective relevant to its own valid regimes. Finally, we shall explore the conservation laws and working regimes relevant to each model.

\quad The fourth chapter is fully consecrated to the study of waves and oscillations in plasmas. Our considerations will begin with the simplest model, the magnetohydrodynamic model, on which we lay the groundwork for examining oscillatory plasma behaviour. We will subsequently extend the work by applying the two-fluid and Vlasov models and comparing their predictions in the relevant regimes. The topics covered include Alfvén waves in the magnetohydrodynamic and the two-fluid description, cold and warm electrostatic plasma waves, as well as elementary cold and warm waves in magnetised plasmas in the two-fluid and the Vlasov description. The latter shall be divided into the studies of perpendicular and parallel propagation to an externally applied magnetic field. The criteria of wave stability and wave damping are also discussed by developing the Landau approach to electrostatic waves using the Vlasov-Poisson model. Our considerations of dynamics in a magnetised environment will later be expanded to treat the kinetics of plasma particles in magnetic fields. 

\quad The fifth chapter is dedicated to the study of the structure of magnetic fields in plasmas. We formally introduce the concepts of magnetic field-lines, magnetic surfaces, and the coordinate systems used to describe them. The main results of our study are magnetic flux coordinate systems, the magnetic-field-line Hamiltonian, and the quasi-symmetry of magnetic fields. These theoretical developments will later serve as a foundation for developing the theory of magnetic confinement.

\quad The sixth chapter firstly formally introduces the reader to the Fokker-Planck collision theory. We revisit our considerations on collisions in plasmas from the first chapter and expand on them. The main theoretical result is the Fokker-Planck equation, which describes the temporal evolution of the plasma particle distribution functions due to collisional processes. Subsequently, we shall develop the theory of neoclassical transport.

\newpage